%%
%% Chapter 4: Experiments and Evaluation (实验与评估)
%%

\chapter{实验与评估}

\section{数据收集}

本研究通过问卷调查收集数据。问卷设计参考了国内外成熟量表，并进行了适当的修改以适应本研究的具体情境。共发放问卷500份，回收有效问卷428份，有效回收率为85.6\%。

\section{描述性统计分析}

对收集到的数据进行描述性统计分析，包括样本的基本特征分布以及各变量的均值、标准差等统计量。分析结果表明，样本具有良好的代表性。

\section{假设检验}

运用统计分析方法对研究假设进行检验。检验结果表明，本研究提出的大部分假设得到了实证支持。具体检验结果如下：

假设1得到支持，变量A对变量B具有显著的正向影响（$\beta = 0.45$，$p < 0.001$）。

假设2得到支持，变量C的中介效应显著（间接效应 = 0.23，95\%置信区间不包含0）。

假设3得到部分支持，变量D的调节效应在低水平时显著，在高水平时不显著。

\section{结果讨论}

对假设检验的结果进行深入讨论，结合理论和已有研究成果，解释研究发现的含义。同时，分析研究中存在的局限性和未来研究方向。
